\appendix

\section{Countably infinite bounded-port extension}
\label{app:infinite}

Let $\mathcal X$ be a separable modal Hilbert space and $K:\mathcal X\rightarrow\mathcal Y$ a bounded total Markov port operator. Let $\mathcal T(t)$ be a passive contraction semigroup satisfying the energy inequality
\begin{equation}
\mathcal T(\tau)\mathcal T^\dagger(\tau)
+\int_0^\tau
\mathcal T(t)K^\dagger K\mathcal T^\dagger(t)\dd t
\le I .
\end{equation}
For a selected bounded input port $K_i$,
\begin{equation}
P_i(\tau)=\int_0^\tau
\mathcal T(t)K_i^\dagger K_i\mathcal T^\dagger(t)\dd t
\end{equation}
obeys
\begin{equation}
0\le P_i(\tau)
\le I-\mathcal T(\tau)\mathcal T^\dagger(\tau)
\le I .
\end{equation}
The monotone strong limit $P_i$ exists and satisfies $0\le P_i\le I$.

Equation~\eqref{eq:kappasum} gives $\Tr(K_g^\dagger K_g)<\infty$ for the retained countable gravitational modal sector, so $K_g$ is Hilbert--Schmidt. For any bounded $L$,
\begin{align}
\int_0^\infty
\|L\mathcal T(t)K_i^\dagger\|_{\HS}^2\dd t
&=\Tr(LP_iL^\dagger)\\
&\le\Tr(L^\dagger L).
\end{align}
The Hilbert--Schmidt operators form a Hilbert space, so vector-valued Plancherel supplies the same frequency-domain weighted $H_2$ cut as in finite dimension. This proof intentionally excludes unbounded boundary-control/observation operators; such systems require separate admissibility analysis \cite{BarasBrockett1975,Opmeer2013}.

\section{STF-sector completeness details}
\label{app:sectorresource}

For any unit-normalized STF tensor $E$, Eq.~\eqref{eq:Einfluence} associates the displacement-space influence field
\begin{equation}
\bm h_E(\bm x)=2E\bm x .
\end{equation}
For an orthonormal tensor set $\{E_\alpha\}$, Bessel's inequality gives
\begin{equation}
\sum_n\frac{1}{\mu_n}
\sum_\alpha|q_n:E_\alpha|^2
\le
\int_V\rho
\sum_\alpha|\bm h_{E_\alpha}(\bm x)|^2\dd^3x .
\end{equation}
Take $\hat R=\hat z$ and transverse coordinates $x,y$. For the two $|m|=2$ tensors,
\begin{equation}
\sum_{\alpha\in |m|=2}|\bm h_{E_\alpha}|^2
=4(x^2+y^2).
\end{equation}
For the two $|m|=1$ tensors,
\begin{equation}
\sum_{\alpha\in |m|=1}|\bm h_{E_\alpha}|^2
=2(x^2+y^2)+4z^2.
\end{equation}
For the single $m=0$ tensor,
\begin{equation}
|\bm h_{E_0}|^2
=\frac23(x^2+y^2)+\frac83z^2.
\end{equation}
Integration gives Eqs.~\eqref{eq:m2resource}--\eqref{eq:m0resource}. Their pointwise sum is $(20/3)r^2$, so the sector decomposition is exactly consistent with the scalar completeness relation.

For a complete displacement basis the three Bessel inequalities become Parseval equalities separately. Consequently the coefficient in the leading $|m|=2$ resource,
\begin{equation}
\sum_n\frac{Q_{2,n}^2}{\mu_n}=4I_{\hat R},
\end{equation}
is sharp at the projection-sum level. If the complete projected resource were also placed in retained modes at $\omega_n=\Omega$, the resource-propagation coefficient $5/4$ in the far-zone source-side inequality would be saturated. This establishes sharpness of the chained mathematical constants within the abstract retained modal class; it is not a claim that an unconstrained homogeneous elastic body realizes that simultaneous modal arrangement.

\section{Exact finite-distance compact TT sectors}
\label{app:finiteTT}

Choose the separation along $\hat z$ and write $\mu=\cos\theta$. Axial rotations diagonalize the normalized compact-STF angular overlap into $m=0$, $|m|=1$, and $|m|=2$ sectors. After azimuthal integration, unit-normalized sector kernels are
\begin{align}
K_2(\mu)&=\frac{5}{32}(1+6\mu^2+\mu^4),\\
K_1(\mu)&=\frac58(1-\mu^2)(1+\mu^2),\\
K_0(\mu)&=\frac{15}{16}(1-\mu^2)^2,
\end{align}
with $\int_{-1}^{1}K_m(\mu)\dd\mu=1$. Their fixed-frequency overlaps are
\begin{equation}
S_m(z)=\int_{-1}^{1}K_m(\mu)e^{iz\mu}\dd\mu,
\qquad z=kR.
\end{equation}
Separating the outgoing $e^{iz}$ part gives
\begin{align}
S_{+,2}(z)&=-\frac{5i}{4z^5}
\left(z^4+2iz^3-3z^2-3iz+3\right)e^{iz},\\
S_{+,1}(z)&=-\frac{5}{2z^5}
\left(z^3+3iz^2-6z-6i\right)e^{iz},\\
S_{+,0}(z)&=\frac{15i}{2z^5}
\left(z^2+3iz-3\right)e^{iz}.
\end{align}
Their squared magnitudes are the three quantities in Eq.~\eqref{eq:finite-sector-eigenvalues}. For $z\ge3$, direct subtraction gives $\eta_2-\eta_1\ge0$ and $\eta_2-\eta_0\ge0$. The separated outgoing component is the retarded source-to-receiver branch used in the wave-zone link. No near-zone channel interpretation of the isolated outgoing component is asserted.

\section{Validation and falsification checks}
\label{app:validation}

Independent numerical regressions were used to challenge, rather than merely illustrate, the analytic inequalities. They are not substitutes for proof.

The regression suite challenges six independent layers. Random passive realizations test the Lyapunov $H_2$ identity and weighted Gramian inequalities. Random mass distributions test the scalar and separate $m=0,1,2$ Parseval/Bessel identities. Random STF tensors and numerical angular integrals test the exact finite-$z$ sector overlaps and singular-value ordering. The assembled theorem is tested against finite-$z$ sector weighting and exact $5/4$ projection-coefficient saturation. Increasing square-summable modal truncations test the countable bounded-port extension, while random contractive reflections test the recurrence resolvent and its phase-aligned scalar saturation.

The analytic development deliberately separates three different issues: propagation variation across the measured band, the compact mass-quadrupole approximation, and the retained near-resonant Markov modal model. Equation~\eqref{eq:finitegeometry} removes the first of these as an unquantified carrier-freezing error. The latter two remain explicit assumptions of the theorem.
